\documentclass{pgnotes}

\title{Automation}

\begin{document}

\maketitle

This class is designed to assist you with CA1.

\section{PSQL input}

\subsection{Shell redirection}

Imagine we have a file \texttt{setup.sql} with the following commands:
\begin{minted}{postgresql}
-- make a table for staff
CREATE TABLE staff ( 
id bigserial primary key,
surname text not null,
firstname text not null,
unique(firstname,surname)
);
-- make a table for projects
CREATE TABLE projects ( 
id bigserial primary key,
description text not null unique,
project_lead bigint references staff
);
\end{minted}

We can run this file without actually going into \texttt{psql} by using the shell redirection operator:

\begin{minted}{bash}
# using default (own name) database
psql < setup.sql

# using other DB (sales) to which have access
psql sales < setup.sql
\end{minted}

\subsection{Psql \textbackslash i}

We can use the the \textbackslash i command within psql to read commands from a file:
\begin{minted}{postgresql}
-- run commands from setup.sql as if typed here
\i setup.sql
\end{minted}

\subsection{Run SQL script from bash via psql}

As displayed by \texttt{psql --help}, we can run commands in a text file \texttt{setup.sql} from bash thus:
\begin{minted}{bash}
psql -f setup.sql 
\end{minted}

\subsection{Run single SQL command from bash}

As displayed by \texttt{psql --help}, we can run a single command straight from bash:
\begin{minted}{bash}
# run select statement and exit
psql -c 'SELECT * FROM staff;' 
\end{minted}

\section{PSQL output}

\subsection{Output shell redirection}

If we start psql using the shell redirection operator we will send all results out to a file:
\begin{minted}{bash}
# default DB
psql > output.txt

# or any DB 
psql sales > output.txt
\end{minted}
No output will be shown in output window. 

\subsection{Psql \textbackslash o}

Alternatively once in psql we can use the \textbackslash o command to send all query results to a file:
\begin{minted}{postgresql}
\o output.txt

-- sql queries
select * from staff;

-- or, inbuilt commands
\dt

-- both will get written to output.txt, not shown on screen
\end{minted}

\section{Table copy}

Table copy operations below will not copy constraints or indices.

\subsection{Copying table definition and data}

\begin{minted}{postgresql}
-- a table w/ columns
TABLE staff;

-- create a copy of table w/ data
CREATE TABLE people AS
TABLE staff;
\end{minted}


\subsection{Copying table definition without data}

\begin{minted}{postgresql}
CREATE TABLE people AS
TABLE staff 
WITH NO DATA;
\end{minted}

\subsection{Copying table definition w/ partial data}

\begin{minted}{postgresql}
CREATE TABLE people AS
SELECT * FROM staff
WHERE office IS NOT NULL;
\end{minted}

\subsection{Table copy from select statement}

Tables can be created and populated from the output of a select statement.
\begin{minted}{postgresql}
-- a query composed of two tables
-- types:      text           int              text                int
SELECT staff.fullname, staff.salary, department.name, department.building
FROM staff
LEFT JOIN department ON department.id=staff.department;

-- would create table w/ types of each col in query
-- and insert data from query
CREATE TABLE staffdirectory AS
SELECT staff.fullname, staff.salary, department.name, department.building
FROM staff
LEFT JOIN department ON department.id=staff.department;

\end{minted}

\section{Generated columns}

\href{https://www.postgresql.org/docs/current/ddl-generated-columns.html}{Generated columns} are the columar equivalent of a view.
Syntax: 
\begin{minted}{postgresql}
CREATE TABLE participant (
    name text primary key,
    height_cm numeric,
    height_in numeric GENERATED ALWAYS AS (height_cm / 2.54) STORED
);
\end{minted}
PostgreSQL only supports STORED generated columns. 


\section{Idempotency}

Issues can arise when the same script is run more than once in succession.
An idempotent operation is one that where it is repeated will have no further effect.

\subsection{Create table if not exists}


\begin{minted}{postgresql}
-- make a table for staff if required
CREATE TABLE IF NOT EXISTS staff ( 
id bigserial primary key,
surname text not null,
firstname text not null,
unique(firstname,surname)
);
\end{minted}


\subsection{Drop table if exists}

\begin{minted}{postgresql}
-- delete table if already there
DROP TABLE IF EXISTS staff;
-- make a table for staff
CREATE TABLE  staff ( 
id bigserial primary key,
surname text not null,
firstname text not null,
unique(firstname,surname)
);
\end{minted}


\end{document}

